\section{Character Paley--Wiener ranges and their actual Gram transport}
\label{sec:character-global-pw}

This section treats nonprincipal primitive Dirichlet characters and
$\psi=\psi_{269,3}$ over $\mathbb Q(\zeta_3)$.  It identifies an
actual finite global range, rather than only the local quotient of
MA-CQ.  The principal rational character is MA-GP-001--008.
The finite-vector construction below also works for every nontrivial
finite-order Hecke character of the latter class-number-one field.
All compact character measures have mass one, as in MA-CH;
Masullo's other Haar normalization requires the rescaling recorded
in (MA-CH.5).

Put $K=\mathbb Q$ or $\mathbb Q(\zeta_3)$, and let $\eta$ be
one of the stated nontrivial primitive characters.  Write
\[
 d=[K:\mathbb Q]\in\{1,2\},\quad
 w_K=|\mathcal O_K^\times|\in\{2,6\},\quad
 Q=|D_K|N\mathfrak f_\eta.
 \tag{MA-CGP.1}
\]
At infinity use parity $e$ when $K=\mathbb Q$, with
$\eta(-1)=(-1)^e$, and radial functions when $K$ is complex.
Let $\mathscr S_\eta$ denote that closed subspace of the
archimedean Schwartz space.  The real Fourier measure is $dx$
and the complex one is $2\,dx\,dy$.  For global Poisson we
explicitly choose the coherent adelic additive character:
at infinity it is $\exp(2\pi i\operatorname{Tr}_{K_\infty/\mathbb R}x)$,
and at finite places it is the negative fractional-trace
character.  Its product is trivial on $K$.  The Fourier
kernel is its complex conjugate.  All Gauss integrals in
the MA-CQ formulas used below are evaluated with this choice.
This does not presume coherence of independently printed
positive-sign local examples in the source.  Conjugating a
finite local additive character reflects that Fourier
coordinate; it preserves its norm support and its value at
zero.  The source's invariant PW support and polar conditions
are therefore unchanged by this explicit normalization.
A Fourier transform changes $\eta$ to
$\bar\eta$; their archimedean parity or radial type is the same.
For an integral ideal put $c(\mathfrak a)=\eta(\mathfrak a)$,
extended by zero on ideals meeting the conductor.

\paragraph{MA-CGP-001 (the primitive theta model and its exact reflected scale).}
For $f\in\mathscr S_\eta$ define
\[
 (\Theta_\eta f)(t)
    =w_K\sqrt t\sum_{\mathfrak a\ne0}
            c(\mathfrak a)f\bigl((N\mathfrak a\,t)^{1/d}\bigr),
                   \qquad t>0.
 \tag{MA-CGP.2}
\]
The values on the positive real ray determine the parity or
radial function.  This is the actual centered, fixed-character
adelic theta image, and
\[
       \Theta_\eta:\mathscr S_\eta
          \xrightarrow{\ \sim\ }\mathcal I_\eta^{\rm cen}
 \tag{MA-CGP.3}
\]
is a topological isomorphism onto the actual smooth ideal
of MT-M22/MT-M24, centered at $1/2$.
There is an explicitly specified phase $\epsilon_\eta^\theta$
of modulus one such that
\[
 \Theta_\eta f(t)
      =\epsilon_\eta^\theta\,
           \Theta_{\bar\eta}(\mathcal F_\infty f)(1/(Qt)).
 \tag{MA-CGP.4}
\]
It is a theta Fourier phase, not an unspecified completed
$L$-function root-number convention.

\begin{proof}
\emph{The actual finite vector and its Fourier transform.}
Use the idele character $\beta$ and local notation
$b_v,\kappa_v,f_v,d_v$ of MA-CQ, with
$\kappa_v=\beta_v^{-1}|_{\mathcal O_v^\times}$.
Take the primitive finite vector
\[
 \varphi_{\eta,v}=
   \begin{cases}
     \mathbf1_{\mathcal O_v},&f_v=0,\\
     \kappa_v\mathbf1_{\mathcal O_v^\times},&f_v>0.
   \end{cases}
 \tag{MA-CGP.5}
\]
A nontrivial character in either field has at least one ramified
finite place: a finite-order character unramified everywhere
would factor through the trivial ideal class group, and its
archimedean connected component is trivial.  At such a place
both the value and the integral of the finite vector vanish.
Consequently $f\otimes\varphi_{\eta,\rm fin}$ satisfies
both adelic polar conditions for every $f\in\mathscr S_\eta$;
no two archimedean moment conditions are imposed.

Choose the positive prime generators over $\mathbb Q$, or
arbitrary prime generators over $\mathbb Q(\zeta_3)$ as in
MA-CQ, and put
\[
 m_v=f_v+d_v,\qquad
 \delta=\prod_v\pi_v^{m_v},\qquad |\delta|_\infty=Q,
\]
\[
 \epsilon_\eta^\theta
      =\prod_{f_v>0}\epsilon_v(\kappa_v)
                                      \prod_v b_v^{m_v}.
 \tag{MA-CGP.6}
\]
The products are finite.  At an unramified place the Fourier
transform of $\mathbf1_{\mathcal O_v}$ is
$q_v^{-d_v/2}\mathbf1_{\pi_v^{-d_v}\mathcal O_v}$.
At a ramified place MA-CQ-001 gives the single shell with
constant $q_v^{-m_v/2}\epsilon_v(\kappa_v)$.
When the separate local dilations $\pi_v^{m_v}$ are replaced
by the one global dilation $\delta$, the other unit components
contribute $\prod_v b_v^{m_v}$, by the principal-idele relation
(MA-CQ.14).  Thus, as an equality of actual finite functions,
\[
 \widehat\varphi_{\eta,\rm fin}(y)
       =Q^{-1/2}\epsilon_\eta^\theta
                                \varphi_{\bar\eta,\rm fin}(\delta y).
 \tag{MA-CGP.7}
\]
The phases in (MA-CGP.6) have modulus one by MA-CQ-001.
Their dual product is $(-1)^e$ in the real case and one in
the complex case; this agrees with $\mathcal F_\infty^2$.
In particular the odd real case has not lost its Fourier sign.

The finite support in (MA-CGP.5) restricts the adelic sum to
integral elements prime to the conductor.  By the principal
character relation its value on such an element $a$ is
$\beta_\infty(a)c((a))$.  Thus it is $c((a))$ on the
positive generator over $\mathbb Q$ and on every generator
in the complex case.  There are $w_K$ generators of each integral ideal.
Their archimedean parity or radial factors compensate the unit
character, so their contributions agree.  This proves that
(MA-CGP.2) is the actual full adelic theta sum on the norm
splitting, with the compact character projected as in MA-CH.

Adelic Poisson summation, with both polar terms zero, gives
$t^{-1}$ times the sum of the Fourier transforms at archimedean
argument $t^{-1/d}a$.  Insert (MA-CGP.7) and replace
$a$ by $\delta^{-1}a$.  At infinity $\delta$ is positive
in the real case; its rotation is immaterial in the complex
radial case.  Multiplying by $\sqrt t$ gives precisely
(MA-CGP.4), including its argument $1/(Qt)$ and constant.
This also gives continuity into exponential Schwartz space:
for large $t$ use Schwartz decay in (MA-CGP.2), and for small
$t$ use (MA-CGP.4).  Ideal counting is $O(T)$ in these two
fields, by interval counting or the rank-two integral lattice;
these estimates also apply after every norm derivative.

\emph{Onto the specified smooth image.}
Here the class-number-one restriction is used essentially.
Every norm-one idele class has a representative in
$K_\infty^1\times\prod_v\mathcal O_v^\times$; its kernel
is the finite global unit group.  Compact character averaging
can therefore be performed on this compact group before
summation.  At infinity it gives exactly the stated parity or
radial Schwartz function.  At a ramified finite place it gives
a finite sum of valuation-shell dilates of
$\kappa_v\mathbf1_{\mathcal O_v^\times}$: compact support
bounds valuations below, and nontrivial unit type together
with local constancy forces vanishing near zero.
At an unramified place the averaged function is constant
on unit shells and constant sufficiently near zero.  It is
a finite linear combination of dilates of
$\mathbf1_{\mathcal O_v}$.  Only finitely many places vary.

A finite-place dilation by the global prime generator
$\pi_v^j$ is moved through the theta sum by reindexing its
principal element.  The other finite unit components supply
the scalar $b_v^j$, and the archimedean factor becomes
$f(\pi_v^j\,\cdot)$, still in $\mathscr S_\eta$.
Thus every projected adelic Schwartz image is a finite linear
combination of images (MA-CGP.2), and hence is one such image.
MT-M24 (M24.7)--(M24.9) identifies this actual image, already
closed, with $\mathcal I_\eta^{\rm cen}$.

Injectivity follows without a spectral assertion.  Put
$F_f(t)=w_K\sqrt t f(t^{1/d})$.
Then (MA-CGP.2) is exactly the ideal sum in (MA-CH.12).
On every compact positive interval its inverse is
\[
 F_f(t)=\sum_{\mathfrak a\ne0}
      \frac{\mu_K(\mathfrak a)c(\mathfrak a)}{\sqrt{N\mathfrak a}}
                   (\Theta_\eta f)((N\mathfrak a)t).
 \tag{MA-CGP.8}
\]
Schwartz decay at infinity and the ideal count make the
series and the substituted double series absolutely convergent,
with all derivatives, on such intervals.  The ideal divisor
sum $\sum_{\mathfrak b\mid\mathfrak a}\mu_K(\mathfrak b)$
is one only for the unit ideal, proving (MA-CGP.8).
Both spaces in (MA-CGP.3) are Fr\'echet, and the target is
closed by MT-M24.  The continuous bijection just proved is
therefore a topological isomorphism by the open mapping theorem.
\end{proof}

\paragraph{MA-CGP-002 (exact finite source support and its Hardy divisor).}
For $\Lambda>0$ let
\[
 \mathscr P_{\eta,\Lambda}
  =\{\Theta_\eta f:f\in\mathscr S_\eta,
          \ \operatorname{supp}\widehat f
                     \subset\{|\xi|_\infty\le\Lambda\}\}.
 \tag{MA-CGP.9}
\]
Then, as an equality of actual smooth spaces,
\[
 \mathscr P_{\eta,\Lambda}
   =\{g\in\mathcal I_\eta^{\rm cen}:
                              \operatorname{supp}g\subset[1/(Q\Lambda),\infty)\}.
 \tag{MA-CGP.10}
\]
This level has an exact source interpretation.  Let $S$ contain
infinity, the conductor and the different, and take $S'=S$
in Masullo Definition~7.3.2.3, pp.~354--355 (thus no additional
unit-shell Euler-deletion transition).  Extend its $S$-local
Schwartz data by $\mathbf1_{\mathcal O_v}$ outside $S$,
apply the actual global character theta map, and use the
invariant Fourier support bound $\prod_{v\in S}|\xi_v|_v\le Y$.
The resulting smooth global level is exactly
\[
                  \mathscr P_{\eta,Y/Q}.
 \tag{MA-CGP.11}
\]
Its Hilbert closure is therefore the actual corresponding
orthogonal range of Definition~11.2.1.3, pp.~488--489,
with this specified choice of finite-level data.

For $0<\alpha<1/2$, polynomial weight $s=0$, and
$c=-\log(Q\Lambda)$, use the equivalent half-line norm
$L^2([c,\infty),e^{2\alpha x}dx)$ and
$(U_cg)(y)=e^{\alpha(c+y)}g(e^{c+y})$.
The Laplace transform is unscaled, with boundary norm
$(2\pi)^{-1}\int|F(it)|^2dt$.
The Hilbert closure $J_{\eta,\Lambda}$ then satisfies
\[
 \begin{aligned}
 \mathcal L U_cJ_{\eta,\Lambda}
       &=B_{\eta,\alpha}H^2(\operatorname{Re}z>0),\\
 Z(B_{\eta,\alpha})&=
   \{\alpha+\tfrac12-\rho_\eta:
                         \operatorname{Re}\rho_\eta<\alpha+\tfrac12\},
 \end{aligned}
 \tag{MA-CGP.12}
\]
with full completed-zero multiplicities.  The inner function is
pure Blaschke: no singular inner factor or positive delay remains.
This does not assume RH for the character.

\begin{proof}
Equation (MA-CGP.4) proves the forward support inclusion in
(MA-CGP.10).  Conversely, apply it and then (MA-CGP.8) in
the dual sector.  If $g$ vanishes below $1/(Q\Lambda)$,
then $\Theta_{\bar\eta}\widehat f(t)=0$ for $t>\Lambda$;
every term of its M\"obius inverse also vanishes there.
Hence $\widehat f(t^{1/d})=0$ for $t>\Lambda$, and parity
or radiality gives the asserted full band support.
The band condition is closed in Schwartz space, so (MA-CGP.3)
also proves that this smooth range is closed.

For (MA-CGP.11), the finite Fourier support of the canonical
vector has maximum norm product $\prod_v q_v^{f_v+d_v}=Q$.
Thus the frequency band $Y/Q$ supplies actual $S$-local
Schwartz data satisfying the source's product support bound
$Y$.  This proves one inclusion, with the stated normalization.
Conversely use global Poisson on arbitrary such source data.
Outside $S$ the Fourier factors are still $\mathbf1_{\mathcal O_v}$,
since all different places are in $S$.
For a nonzero principal element $a$ contributing to the reflected
sum, $|a|_v\le1$ outside $S$, and the product formula gives
$|a|_S\ge1$.  At norm $t<1/Y$, its reflected $S$-argument
has norm at least $1/t>Y$, so every summand is zero.
Compact character averaging preserves this statement: the
class-number-one compact representatives above have norm one
in $S$ and unit components outside it.
The actual projected image therefore lies in the smooth ideal
and is supported in $[1/Y,\infty)$.  Equation (MA-CGP.10)
gives the reverse inclusion.  Smooth closure preserves the
same support condition, proving (MA-CGP.11) at the source's
closed-image level as well.

We give the Hardy normalization and common divisor explicitly.
Right translation on the support half-line preserves the smooth
ideal and (MA-CGP.10).  Under $U_c$ it is a scalar times the
unilateral translation.  The Laplace image of the closure is
therefore a nonzero invariant subspace $BH^2$, by the same
vectorwise resolvent average and Beurling theorem proved in
MA-GP-003.
For a frequency test $u=\widehat f$ supported in a positive
annulus inside the band, put $U(r)=u(r^{1/d})$ and
$s_z=z+1/2-\alpha$.  Direct Mellin integration, first where
it converges absolutely, gives
\[
 (\mathcal L U_c\Theta_\eta f)(z)
   =w_K\epsilon_\eta^\theta e^{zc}Q^{z-\alpha}
          L_K(s_z,\bar\eta)
                   \int_0^\Lambda U(r)r^{s_z-1}\,dr.
 \tag{MA-CGP.13}
\]
Analytic continuation gives the equality on the whole right
half-plane.  There $\operatorname{Re}s_z>1/2-\alpha>0$,
so the gamma factors have no poles, and the nontrivial finite
$L$-function is entire.  No polar cancellation or moment-zero
constraint is required.  At every prescribed $s_z$ the final
integral can be made nonzero by an annular test.  Thus the
common zeros are exactly those of $L_K(s_z,\bar\eta)$,
to their exact orders.  The completed functional equation
$\Lambda_\eta(s)=W_\eta\Lambda_{\bar\eta}(1-s)$ changes
their indexing to precisely (MA-CGP.12).  This explains why
the dual character in (MA-CGP.13) is essential.

Every core output is flat at $c$ and decreases faster than
all exponentials at $+\infty$.  Its Laplace transform is
entire and, by repeated integration by parts, decreases faster
than every power uniformly on the closed right half-plane.
The disk image used in MA-GP-003 is consequently bounded
and analytically extends across every boundary point other
than infinity.  The same elementary singular-measure argument
there excludes finite boundary atoms and all finite singular
support; only a delay $e^{-az}$ at infinity could remain.
To exclude it, choose $u$ positive on the positive ray arbitrarily
close to the band boundary, flat there and zero near the origin,
with the required parity or radial extension.  For $t$ just
above $1/(Q\Lambda)$ only the unit ideal contributes to the
reflected formula (MA-CGP.4).  Its value is a nonzero constant
phase times that positive test.  The range has no positive
delay.  This proves the pure Blaschke assertion, as well as
Blaschke summability through its actual Hardy realization.
\end{proof}

The restriction $S'=S$ in (MA-CGP.11) is substantive.  If
additional unit-shell transition factors are inserted at
unramified places, their finite Euler multipliers and changed
levels must be retained; this theorem does not identify those
Euler-deleted ranges with the primitive range.  The fixed-base
levels just proved already form a dense increasing band exhaustion of
the full primitive character image.  The absolute hypercube
levels contain every one of its canonical tests once their
level exceeds the finite coordinate bounds and the
archimedean band, so their union is dense as well.

\paragraph{MA-CGP-003 (all-weight condensation and a concrete character Gram kernel).}
For every $\alpha>0$ and $s\ge0$, the orthogonal projections
$Q_{\eta,\Lambda}$ onto the closures of (MA-CGP.9) increase
strongly to the actual geometric projection $P_{G,\eta}$.
For $A=M_\chi C_h$ with compact smooth norm tests,
\[
       \|AQ_{\eta,\Lambda}-AP_{G,\eta}\|_1\longrightarrow0.
 \tag{MA-CGP.14}
\]
The actual half-line Gram formula of MA-GP-004 applies to
these ranges at every polynomial weight; for $s=0$ its
inner multiplier is the explicitly identified $B_{\eta,\alpha}$.

There is also an explicit frequency Gram transport.  Write
\[
 a_n=\sum_{N\mathfrak a=n}\overline{c(\mathfrak a)},\qquad
 \gamma_K=\begin{cases}2,&K=\mathbb Q,\\2\pi,&K=\mathbb Q(\zeta_3),\end{cases}
\]
\[
 \mathcal P_{\eta,\Lambda}=L^2((0,\Lambda),\gamma_K\,dy).
 \tag{MA-CGP.15}
\]
This is precisely the auxiliary ordinary archimedean band
norm on $U(y)=u(y^{1/d})$; no weighted local norm is implicit.
On the dense smooth parity or radial band core, including
functions nonzero at the origin when the type permits them,
Poisson gives
\[
 (T_{\eta,\Lambda}U)(t)
   =\frac{w_K\epsilon_\eta^\theta}{\sqrt{Qt}}
                      \sum_{n\ge1}a_n U(n/(Qt)).
 \tag{MA-CGP.16}
\]
It is closable and injective after closure, and its range
closure is $J_{\eta,\Lambda}$.
For the norm (MA-GP.52) and $A=M_\chi C_h$, its core form
has the bounded trace-class extension
\[
 \begin{aligned}
 (K_{\eta,\chi,h}U)(y)
   ={}&\frac{w_K^2}{\gamma_K}
      \sum_{m,n\ge1}\frac{\overline{a_m}a_n}{m}
       \chi\!\left(\frac{m}{Qy}\right)
       w_{\alpha,s}\!\left(\log\frac{m}{Qy}\right)\\
      &\hspace{12mm}\times
       \int h(r)r^{1/2}U(nry/m)\,\frac{dr}{r}.
 \end{aligned}
 \tag{MA-CGP.17}
\]
For $\operatorname{supp}\chi\subset[a,b]$ and
$\operatorname{supp}h\subset[r_-,r_+]$, only
$m\le Q\Lambda b$ and $n\le Q\Lambda b/r_-$ contribute.
Let $G_\eta=\overline T_{\eta,\Lambda}^{\,*}\overline T_{\eta,\Lambda}$.
Then the actual fixed-level trace is
\[
 \operatorname{Tr}_H(AQ_{\eta,\Lambda})
    =\lim_{\varepsilon\downarrow0}
       \operatorname{Tr}_{\mathcal P_{\eta,\Lambda}}
       \bigl((G_\eta+\varepsilon)^{-1/2}
                K_{\eta,\chi,h}(G_\eta+\varepsilon)^{-1/2}\bigr),
 \tag{MA-CGP.18}
\]
with the same two actual trace-norm limits as (MA-GP.68).
It is a Hilbert trace transport; representation intertwining
or an LF lift of its Gram resolvents is not asserted.

\begin{proof}
Smooth frequency cutoffs preserving parity or radiality
approximate every $f\in\mathscr S_\eta$ in Schwartz topology.
Both adelic polar conditions remain automatic from the fixed
nontrivial finite vector.  Equations (MA-CGP.3) and (MA-CGP.9)
give density of the increasing union in the actual smooth
ideal and in its weighted closure.  Strong projection convergence
follows.  Since $A$ itself is trace class by MA-LIFT-003,
finite-rank approximation and self-adjoint strong convergence
prove (MA-CGP.14), for all the stated weights.
The equivalent half-line norm has the same ratio
$R(x)=e^{2\alpha|x|-2\alpha x}$ as (MA-GP.40)--(MA-GP.44);
its positive bounded inverse proof applies without change.

The source norm in (MA-CGP.15) follows directly from
$2\int_0^\infty|u(r)|^2dr$ on a real parity sector and
$4\pi\int_0^\infty|u(r)|^2r\,dr$ on the complex radial
sector, substituting $y=r^d$.
The dense core consists of actual smooth archimedean functions
with frequency support in the open band.  It need not vanish
near the origin.  Dilation into a slightly smaller band
shows that this core gives the same range closure as the
closed-band definition (MA-CGP.9).
Every compact output localization of (MA-CGP.16) has only
finitely many terms and is bounded on this auxiliary source,
by Cauchy--Schwarz and the substitutions $y=n/(Qt)$.
This proves closability exactly as in MA-GP-005.
The finite inverse coefficients
$b_n=\sum_{N\mathfrak a=n}\mu_K(\mathfrak a)
                              \overline{c(\mathfrak a)}$
satisfy $\sum_{d\mid n}a_db_{n/d}=\mathbf1_{n=1}$.
Applied to the locally finite Poisson sum on a graph-domain
vector, they prove injectivity of the closure, with the same
null-set justification as (MA-GP.72).

For the form calculation, the squared scalar in (MA-CGP.16)
is $w_K^2/Q$.  In the pairing with the $m$-term put
$y=m/(Qt)$; the measure $t^{-2}dt$ supplies $Q\,dy/m$.
Dividing by the source density $\gamma_K\,dy$ gives
exactly $w_K^2/(\gamma_K m)$ and the conjugate coefficient
$\overline{a_m}a_n$ in (MA-CGP.17).
After $z=nry/m$, each kernel is a bounded weight multiplier
times
\[
       \frac{w_K^2\overline{a_m}a_n}
                    {\gamma_K\sqrt{mn}\sqrt{yz}}
                              h\!\left(\frac{mz}{ny}\right).
 \tag{MA-CGP.19}
\]
Its support is away from both coordinate origins.  The finite
smooth-kernel trace-class argument of MA-GP-008 applies,
without a zero-mean projection: no such condition belongs
to this nontrivial finite-vector source.
Graph-core approximation extends the form identity to
$\operatorname{Dom}\overline T$ in both variables.
Spectral calculus then maps the whole source into that domain
under $(G_\eta+\varepsilon)^{-1/2}$ and gives a contraction
$F_\varepsilon=\overline T(G_\eta+\varepsilon)^{-1/2}$.
It and its adjoint converge strongly to the actual polar
isometry and its adjoint.  Applying the proved trace-ideal
continuity and cross-space cyclicity from MA-GP-008 gives
(MA-CGP.18), including both trace-norm limits.  These are
actual bounded maps, not formal products with an unspecified
adjoint domain.
\end{proof}

\paragraph{The two requested character families and the remaining asymptotic.}
For primitive Dirichlet $\chi$ the scale is $Q=q$, and the
Hardy zeros in (MA-CGP.12) are the zeros of its own primitive
completed function, with parity retained throughout Fourier
transformation.  For $\psi_{269,3}$ the scale is
$Q=3\cdot269^2$, with $\delta=(1-\zeta_3)269$ and
$b_3=b_{269}=1$.  Thus (MA-CGP.6) uses the Gauss phase of
$\psi^{-1}$ on the ramified unit shell; the quotient cutoff
in MA-CQ-004 uses the dual phase, as it should.
In both cases the source invariant level $Y$ has the same
support edge $-\log Y$, after the exact primitive band
conversion $\Lambda=Y/Q$.

The fixed-character global orthogonal range, its common Hardy
divisor, its bounded half-line Gram and its explicit local cyclic
operator are now specified.  This reduces task~1(d)'s remaining
global/local issue to an asymptotic comparison of that actual
Gram-corrected operator with the separately constructed local
character cutoff of MA-CQ.  Its Gauss and conductor factors
cannot be dropped.  Nothing here changes the all-weight
retained multiplicities of MA-WP or supplies an independent
narrow-strip return.  Dirichlet RH is the stated scope of $A$;
the cubic character is outside that scope and has the distinct
single-character Hecke RH question already identified in MA-CH.

\paragraph{Source and Grothendieck scope.}
MT-M22 (M22.1)--(M22.10) and MT-M24, especially
(M24.2), (M24.6)--(M24.9), were read completely for the actual
character, Poisson and closed-image identifications used here.
The mixed support definitions and proof were read at Masullo
pp.~350--357, especially Definition~7.3.2.3 and
Proposition~7.3.2.4; the Hilbert projection definition was
read at pp.~487--489.  The exact $S'=S$ specialization is
stated above instead of discarding the source's extra Euler
transition factors.  The finite Fourier calculations are the
independently checked MA-CQ-001--002.  The Hardy argument
uses the Beurling/Burnol source scope already given in
MA-GP-003, with every new character and conductor factor
derived here.

The rank-one finite-monodromy character and its induced Artin
object are those of MA-CH and MA-CQ.  The finite-length
$\mathbb C[z]$ zero modules retain the full common-zero
multiplicities before Hilbert truncation; a Gram normalization
changes the analytic metric, not that algebraic length.
The actual trace-class cyclic map is the ordinary Hilbert
counterpart of the inherited cote~1 trace/determinant programme,
BR-TR-01/02, scans~079--080 and~111--112.  Those archive
locators are inherited, not newly inspected here.  No
arithmetic positivity, half-integral Tate object, or new
$\ell$-adic weight realization is claimed.
